\documentclass[a4paper]{report}

\input{config}

% Document title and author
\title{\huge \textbf{\textcolor{myColor}{Template\_latex}}\\}
\author{Tarek Naja Sara Vaccaro - 2026}
\date{}

\begin{document}

\thispagestyle{empty} % rimuove intestazione/piede pagina

% decorative thin left band (1.2cm) drawn in overlay so it doesn't change layout
\begin{tikzpicture}[remember picture,overlay]
    % extend a couple millimetres below the page bottom to avoid any tiny gap
    \fill[ArancioneFascia] (current page.north west) rectangle ($(current page.south west)+(1.4cm,0)$);
\end{tikzpicture}

\begin{center}
    % Titolo Università
\vphantom{space} \\[0.5cm]
{\Huge \textbf{\textcolor{myColor}{PAWS}}}\\[0.3cm]

    % Corso di Laurea
    {\normalsize Programmazione Mobile - Applicazione in Kotlin e Flutter}\\[2cm]
    
    % Logo
    
\includegraphics[width=0.35\textwidth]{univpm_logo.png}\\[2cm]
    
    % Titolo progetto
    {\large \textbf{\textcolor{myColor}{UNIVERSITÀ POLITECNICA DELLE MARCHE}}}\\ [0.3cm]
    % Corso di Laurea
    {\normalsize Corso di Laurea Triennale in Ingegneria Informatica e dell’Automazione}\\[2cm]
    
    % Relatori e Tesina
    \begin{minipage}{0.45\textwidth}
        \begin{flushleft}
        \textbf{\textcolor{myColor}{RELATORI:}}\\[0.2cm]
         Emanuele Storti
        \end{flushleft}
    \end{minipage}
    \hfill
    \begin{minipage}{0.45\textwidth}
        \begin{flushright}
        \textbf{\textcolor{myColor}{TESINA DI:}}\\[0.2cm]
        Tarek Naja
        \\
        Sara Vaccaro
        \end{flushright}
    \end{minipage}\\[5.5cm]
    
    % Anno accademico
    {\normalsize A.A 2025/2026}
\end{center}

% Include the table of contents
\tableofcontents
% Remove page number from table of contents page
\thispagestyle{empty}

\setcounter{page}{0}

\clearpage
\newchapter{Introduzione}
\newsection{Descrizione in linguaggio naturale}
Il progetto consiste nello sviluppo di \textbf{Paws}, un'applicazione mobile pensata per amanti degli animali privati e rifugi. L'obiettivo dell'applicazione è offrire uno strumento per facilitare l'adozione e la condivisione di annunci relativi a cuccioli ed animali domestici, tenendo traccia delle schede di dettaglio dei singoli animali e consentendo il contatto diretto tra gli utenti.

\noindent Il sistema prevede due macro-tipologie di utenti:
\begin{itemize}
   \item \textbf{Utenti privati:} individui che utilizzano l'applicazione per consultare gli annunci disponibili, salvare gli animali preferiti nella propria sezione dedicata, pubblicare nuovi annunci di adozione, filtrare la ricerca per razza, età, sesso e tipologia di utente, contattare direttamente i proprietari tramite chiamata telefonica, nonché cercare altri utenti e rifugi per animali (animal shelters) per visionarne i profili e i relativi post.
   \item \textbf{Rifugi e canili (Animal shelters):} strutture ed organizzazioni dedicate che dispongono di funzionalità per la pubblicazione e gestione di annunci di adozione su più ampia scala, promuovendo la visibilità degli animali in cerca di una casa e gestendo il proprio profilo informativo.
\end{itemize}

\section*{Sintesi dell'Approccio}
Per la realizzazione del sistema, lo sviluppo dell'applicazione mobile è stato suddiviso in due percorsi complementari:
\begin{itemize}
   \item \textbf{Applicazione Android nativa:} sviluppata in Kotlin sfruttando l'architettura nativa Android e la libreria di componenti UI, garantendo elevate prestazioni, gesture fluide (swipe-to-favorite) ed integrazione sia con i servizi di sistema (come l'avvio delle chiamate) sia con servizi di backend esterni. Nello specifico, sono stati utilizzati Firebase per l'autenticazione e la memorizzazione delle credenziali degli utenti e Supabase per il caricamento delle immagini e l'archiviazione dei dati completi relativi ai cuccioli.
   \item \textbf{Applicazione clone cross-platform:} realizzata in Flutter per consentire la portabilità del codice a partire da un'unica codebase, replicando fedelmente le funzionalità, l'identità visiva e l'interfaccia dell'applicazione nativa.
\end{itemize}

\noindent Questo approccio parallelo consente di analizzare e confrontare i vantaggi dello sviluppo nativo in termini di integrazione con il sistema operativo e risposta delle gesture rispetto ai benefici legati alla velocità di sviluppo e riusabilità offerti dalle moderne tecnologie cross-platform


\newchapter{Glossario dei Termini}


Il seguente glossario definisce in modo univoco i termini di business e tecnologici utilizzati nell’ambito del progetto per evitare ambiguità durante le fasi di analisi, progettazione e sviluppo.


\vspace{0.5cm} % Spazio opzionale per separare il testo dalla tabella


\begin{longtable}{ >{\raggedright\arraybackslash}p{1.5cm} >{\raggedright\arraybackslash}p{5.5cm} >{\raggedright\arraybackslash}p{2cm} >{\raggedright\arraybackslash}p{2.97cm}}
\toprule
\textbf{TERMINE} & \textbf{DESCRIZIONE} & \textbf{TIPO} & \textbf{SINONIMI} \\
\midrule
\endhead
\bottomrule
\endfoot
\endlastfoot
Utente Privato & Persona fisica registrata che utilizza l'app per cercare cuccioli, salvare preferiti e pubblicare annunci di adozione & BUSINESS & Privato, User \\ \addlinespace
Animal Shelter & Struttura d'accoglienza o associazione no-profit che utilizza l’app per cercare cuccioli, salvare preferiti e pubblicare annunci di adozione & BUSINESS & Rifugio, Ente di adozione \\ \addlinespace
Cucciolo & Entità principale dell'applicazione che rappresenta un animale pubblicato con foto, nome, età, sesso, tipo e descrizione & BUSINESS & Animale, Post, Pet, Puppy \\ \addlinespace
Scheda Dettaglio & Schermata informativa dedicata ad un singolo cucciolo contenente tutti i dettagli dell'animale e i contatti del proprietario & BUSINESS & Dettagli Cucciolo, Puppy Details \\ \addlinespace
Preferiti & Sezione dell'applicazione in cui l'utente raccoglie e salva i cuccioli a cui è interessato & BUSINESS & Saved, Mi Piace, Favorites \\ \addlinespace
Notifica & Messaggio automatico inviato al proprietario di un annuncio quando un altro utente aggiunge/rimuove il suo cucciolo ai preferiti & BUSINESS & Avviso, Notification \\ \addlinespace
Filtro di Ricerca & Strumento per restringere l'elenco degli animali visualizzati in base a parametri quali età, sesso, tipo di animale e tipo di utente & BUSINESS & Criterio di Selezione, Search Filter \\ \addlinespace
Account & Insieme delle credenziali di accesso (email e password) e dei dati personali che identificano un utente nel sistema & TECNICO & Profilo \\ \addlinespace
Swipe to Favorite & Gesto di scorrimento orizzontale della scheda verso destra che consente di salvare un cucciolo nei preferiti & TECNICO & Gesture Swipe, Scorrimento Destra \\ \addlinespace
Firebase & Servizio cloud utilizzato per la gestione della registrazione, dell'autenticazione sicura e della sessione degli utenti & TECNICO & Auth, Gestione Accessi \\ \addlinespace
Cloud Firestore & Database NoSQL in tempo reale utilizzato per memorizzare i dati degli utenti, dei post, dei preferiti e delle notifiche & TECNICO & DB NoSQL, Firestore \\ \addlinespace
Supabase & Servizio di Object Storage utilizzato per l'hosting e l'erogazione performante delle immagini del profilo e delle foto dei cuccioli & TECNICO & Storage Cloud \\ \addlinespace
\end{longtable}

\newchapter{Progettazione e sviluppo applicazione}
Nella presente sezione si analizzano in modo approfondito i requisiti del sistema, procedendo alla loro suddivisione secondo le seguenti categorie:
\begin{itemize}
   \item Requisiti funzionali, che definiscono le specifiche funzionalità che il sistema è tenuto a fornire.
   \item Requisiti non funzionali, che individuano i vincoli e le qualità che il sistema deve soddisfare, quali ad esempio l'usabilità, la sicurezza e le prestazioni.
\end{itemize}
\newsection{Analisi dei requisiti}
\newsubsection{Requisiti Funzionali}
\subsection*{Area: Gestione Utente}
\begin{itemize}
    \item \textbf{RF1: Registrazione utente:} il sistema deve permettere la creazione di un account inserendo email, password, nome, cognome, città, provincia, CAP, numero di telefono, username e la selezione del tipo di account (Utente Privato o Animal Shelter).
    \item \textbf{RF2: Accesso utente:} il sistema deve consentire agli utenti registrati di accedere all'applicazione inserendo le proprie credenziali.
    \item \textbf{RF3: Modifica utente:} il sistema deve permettere all'utente di modificare i propri dati personali, aggiornare la foto profilo, richiedere il reset della password via email e modificare la tipologia di account.
    \item \textbf{RF4: Logout:} il sistema deve consentire all'utente di effettuare un logout sicuro dalla sessione attiva.
    \item \textbf{RF5: Elimina utente:} il sistema deve permettere all'utente di eliminare definitivamente il proprio account.
    \item \textbf{RF6: Visualizza utente:} il sistema deve permettere all'utente di visualizzare la pagina di profilo di un altro utente consultando le informazioni pubbliche e i post pubblicati.
\end{itemize}
\subsection*{Area: Gestione Cucciolo}
\begin{itemize}
    \item \textbf{RF7: Crea cucciolo:} il sistema deve permettere all'utente di caricare un nuovo annuncio specificando nome del cucciolo, tipo di animale (cane, gatto, uccello), sesso (maschio, femmina), età, didascalia descrittiva e foto dell'animale.
    \item \textbf{RF8: Modifica cucciolo:} il sistema deve permettere all'autore dell'annuncio di modificare le informazioni inserite.
    \item \textbf{RF9: Elimina cucciolo:} il sistema deve permettere all'autore dell'annuncio di eliminare definitivamente il post e la relativa immagine memorizzata nel cloud.
    \item \textbf{RF10: Visualizza post e scheda dettaglio:} il sistema deve mostrare i post dei cuccioli pubblicati e consentire la visualizzazione delle informazioni complete dell'animale e del proprietario.
\end{itemize}
\subsection*{Area: Gestione nterazioni}
\begin{itemize}
    \item \textbf{RF11: Aggiungi ai preferiti:} il sistema deve permettere di aggiungere un cucciolo ai preferiti effettuando uno swipe verso destra dalla home o premendo sul cuore dal profilo di un Utente o Animal Shelter.
    \item \textbf{RF12: Rimuovi dai preferiti:} il sistema deve permettere di rimuovere un cucciolo dai preferiti premendo sul cuore nella lista dei preferiti.
    \item \textbf{RF13: Calcolo dei mi piace:} il sistema deve calcolare e mostrare in tempo reale il numero totale di salvataggi ricevuti per ciascun post.
    \item \textbf{RF14: Chiamata telefonica:} il sistema deve consentire di avviare una chiamata verso il numero di telefono del proprietario.
\end{itemize}
\subsection*{Area: Gestione Navigazione}
\begin{itemize}
    \item \textbf{RF15: Cerca:} il sistema deve permettere di cercare utenti o cuccioli inserendo il testo nella barra di ricerca.
    \item \textbf{RF16: Filtra:} il sistema deve consentire il filtraggio degli annunci in base a età, sesso, tipo di animale e tipo di utente.
\end{itemize}
\subsection*{Area: Gestione Notifiche}
\begin{itemize}
    \item \textbf{RF17: Genera notifica:} il sistema deve generare una notifica destinata al proprietario dell'annuncio quando un altro utente aggiunge/rimuove il suo cucciolo ai/dai preferiti.
    \item \textbf{RF18: Visualizza notifica:} il sistema deve mostrare all'utente un elenco delle notifiche ricevute.
    \item  \textbf{RF19: Elimina notifica:} il sistema deve permettere all'utente l'eliminazione delle notifiche.
\end{itemize}

\newsubsection{Requisiti Non Funzionali: Kotlin (App Nativa)}
\begin{itemize}
    \item \textbf{RNF1-K: Tecnologia di Sviluppo e SDK:} L’applicazione nativa è sviluppata in Kotlin con Android SDK, strutturata mediante Activity e Fragment per la navigazione tra le varie sezioni della schermata principale.
    \item \textbf{RNF2-K: Layout ed Interfaccia Grafica:} L’interfaccia utente è realizzata interamente mediante Layout XML, utilizzando i componenti Material Design e la libreria ViewBinding per il collegamento sicuro tra codice Kotlin e layout.
    \item \textbf{RNF3-K: Autenticazione e Database Cloud:} La gestione della sessione utente e la persistenza dei dati di dominio (Post, Profili Utente, Notifiche, Preferiti) sono gestite tramite Firebase Auth per l'autenticazione ed a Cloud Firestore per la lettura/scrittura in tempo reale dei dati.
    \item \textbf{RNF4-K: Cloud Storage per i Media:} L'upload e la gestione dello storage remoto delle immagini di profili e cuccioli sono realizzati tramite l'integrazione diretta con Supabase Storage.
    \item \textbf{RNF5-K: Rendering e Caching delle Immagini:} Il caricamento asincrono, il rendering ed il caching locale delle immagini da URL remoto nelle ImageView dell'interfaccia utente sono affidati alla libreria Glide.
    \item \textbf{RNF6-K: Visualizzazione Dati e RecyclerView:} L'applicazione organizza la presentazione delle liste e delle griglie di dati (post, cuccioli, notifiche, preferiti) mediante RecyclerView e Adapter dedicati (FeedAdapter, GridPostAdapter, FavoritesAdapter, NotificationAdapter, SearchAdapter).
\end{itemize}

\newsubsection{Requisiti Non Funzionali: Flutter (App Cross-Platform)}
\begin{itemize}
    \item \textbf{RNF1-F: Tecnologia di Sviluppo Cross-Platform:} L’applicazione cross-platform è sviluppata in Dart con il framework Flutter (SDK \textasciicircum 3.12.2), garantendo la compatibilità sia per dispositivi Android sia per iOS con un'unica codebase.
    \item \textbf{RNF2-F: Servizi e Modelli Dati:} La gestione dei dati di dominio e della logica di backend è strutturata mediante Service dedicati (FirebaseService, SupabaseService) e Modelli di dati fortemente tipizzati (UserModel, PuppyPostModel, NotificationModel).
    \item \textbf{RNF3-F: Autenticazione e Cloud Firestore:} La gestione dell'autenticazione utente e le operazioni di lettura/scrittura sulle collezioni cloud (Cuccioli, Post, Preferiti, Notifiche, Utenti) sono affidate ai pacchetti ufficiali firebase\_auth e cloud\_firestore.
    \item \textbf{RNF4-F: Cloud Storage su Supabase:} Il caricamento e l'hosting remoto dei file multimediali (immagini dei cuccioli nel bucket puppy-images ed avatar utenti nel bucket profile-images) sono realizzati tramite il pacchetto supabase\_flutter.
    \item \textbf{RNF5-F: Gestione dell'Interfaccia e dello Stato:} L'interfaccia utente e l'aggiornamento dello stato reattivo sono realizzati mediante i componenti nativi StatefulWidget e StatelessWidget di Flutter, con aggiornamento dinamico gestito tramite setState() e l'invocazione diretta dei Service.
    \item \textbf{RNF6-F: Design System e Caching Immagini:} L'applicazione impiega la classe centralizzata AppTheme (lib/theme/app\_theme.dart) per la definizione di colori (AppColors), stili ed elementi grafici, e la libreria cached\_network\_image per la visualizzazione ottimizzata delle immagini di rete con placeholder di caricamento.
\end{itemize}


\clearpage
\newsection{Descrizione dei Casi d'Uso}
Di seguito sono descritti nel dettaglio i principali casi d’uso per l’applicazione.
\newsubsection{Area: Gestione Utente}
Questa area raggruppa tutti i casi d'uso relativi all'autenticazione e all'amministrazione del profilo personale. Definisce i processi tramite gli utenti possono registrarsi, accedere in modo sicuro tramite le proprie credenziali, aggiornare le informazioni anagrafiche e gestire l'intero ciclo di vita del proprio account, inclusa la disconnessione e l'eliminazione definitiva dei dati.
\vspace{1cm}
\begin{figure}[ht!]
    \centering
    \includegraphics[width=0.8\textwidth]{GestioneUtente.png}
    \caption{Gestione Utente}
\end{figure}

% UC-1: Registrazione Utente
\begin{table}[H]
\vspace*{-0cm}
\renewcommand{\arraystretch}{1.9}
\begin{tabular}{|p{3.9cm}|p{9.9cm}|}
\hline
\multicolumn{2}{|c|}{\textbf{Caso d’uso: RegistrazioneUtente}} \\ \hline
\textbf{ID} & UC-1 \\ \hline
\textbf{Breve descrizione} & Un Utente crea un nuovo account per accedere all'applicazione. \\ \hline
\textbf{Attori primari} & Utente non autenticato \\ \hline
\textbf{Attori secondari} & Sistema di Autenticazione (Firebase Auth), Database (Cloud Firestore) \\ \hline
\textbf{Precondizioni} & \begin{enumerate}[leftmargin=14pt,label=\arabic*.,labelsep=0.5em,topsep=0pt,partopsep=0pt,parsep=0pt,itemsep=0pt]
    \item L'applicazione è avviata e l'Utente si trova nella schermata di registrazione.
\end{enumerate} \\ \hline
\textbf{Sequenza degli eventi principale} &
\begin{enumerate}[leftmargin=14pt,label=\arabic*.,labelsep=0.5em,topsep=0pt,partopsep=0pt,parsep=0pt,itemsep=0pt]
   \item L'Utente inserisce i propri dati: email, password, nome, cognome, città, provincia, CAP, numero di telefono e username.
   \item L'Utente seleziona la tipologia di account (Utente Privato o Animal Shelter).
   \item L'Utente convalida la registrazione premendo il pulsante dedicato.
   \item Il sistema valida i dati inseriti.
   \item Il sistema registra il nuovo account su Firebase Auth e ne salva le informazioni su Cloud Firestore.
   \item L'Utente viene reindirizzato alla schermata principale e risulta autenticato.
\end{enumerate}\\ \hline
\textbf{Postcondizioni} &
\begin{enumerate}[leftmargin=14pt,label=\arabic*.,labelsep=0.5em,topsep=0pt,partopsep=0pt,parsep=0pt,itemsep=0pt]
\item L'account viene creato ed è abilitato all'uso, e l'Utente risulta autenticato nel sistema.
\end{enumerate} \\ \hline
\textbf{Sequenza degli eventi alternativa} & \begin{enumerate}[leftmargin=14pt,label=\arabic*.,labelsep=0.5em,topsep=0pt,partopsep=0pt,parsep=0pt,itemsep=0pt]
   \item \texttt{If} i dati non sono validi o l'email risulta già registrata, allora
   \begin{enumerate}[label=\arabic{enumi}.\arabic*.,leftmargin=22pt,labelsep=0.5em,topsep=0pt,partopsep=0pt,parsep=0pt,itemsep=0pt]
       \item Il sistema mostra un messaggio d'errore.
       \item Il sistema invita l'Utente a correggere i campi errati o a inserire un indirizzo differente.
   \end{enumerate}
\end{enumerate}\\ \hline
\end{tabular}
\end{table}

% UC-2: Accesso Utente (Login)
\begin{table}[H]
\vspace*{-0cm}
\renewcommand{\arraystretch}{1.9}
\begin{tabular}{|p{3.9cm}|p{9.9cm}|}
\hline
\multicolumn{2}{|c|}{\textbf{Caso d’uso: AccessoUtente}} \\ \hline
\textbf{ID} & UC-2 \\ \hline
\textbf{Breve descrizione} & Un Utente registrato effettua l'accesso all'applicazione inserendo le proprie credenziali. \\ \hline
\textbf{Attori primari} & Utente registrato \\ \hline
\textbf{Attori secondari} & Sistema di Autenticazione (Firebase Auth) \\ \hline
\textbf{Precondizioni} & \begin{enumerate}[leftmargin=14pt,label=\arabic*.,labelsep=0.5em,topsep=0pt,partopsep=0pt,parsep=0pt,itemsep=0pt]
    \item L'Utente possiede un account valido e si trova nella schermata di login.
\end{enumerate} \\ \hline
\textbf{Sequenza degli eventi principale} &
\begin{enumerate}[leftmargin=14pt,label=\arabic*.,labelsep=0.5em,topsep=0pt,partopsep=0pt,parsep=0pt,itemsep=0pt]
   \item L'Utente inserisce la propria email e password.
   \item L'Utente preme il pulsante di accesso.
   \item Il sistema verifica la correttezza delle credenziali.
   \item Il sistema autentica l'Utente e lo reindirizza alla schermata principale.
\end{enumerate}\\ \hline
\textbf{Postcondizioni} &
\begin{enumerate}[leftmargin=14pt,label=\arabic*.,labelsep=0.5em,topsep=0pt,partopsep=0pt,parsep=0pt,itemsep=0pt]
\item L'Utente è autenticato con successo e può accedere alle funzionalità riservate del sistema.
\end{enumerate} \\ \hline
\textbf{Sequenza degli eventi alternativa} & \begin{enumerate}[leftmargin=14pt,label=\arabic*.,labelsep=0.5em,topsep=0pt,partopsep=0pt,parsep=0pt,itemsep=0pt]
   \item \texttt{If} le credenziali inserite risultano errate o inesistenti, allora
   \begin{enumerate}[label=\arabic{enumi}.\arabic*.,leftmargin=22pt,labelsep=0.5em,topsep=0pt,partopsep=0pt,parsep=0pt,itemsep=0pt]
       \item Il sistema mostra un messaggio d'errore per credenziali non valide.
       \item Il modulo viene pulito o l'Utente è invitato a riprovare l'inserimento.
   \end{enumerate}
\end{enumerate}\\ \hline
\end{tabular}
\end{table}

% UC-3: Modifica Utente
\begin{table}[H]
\vspace*{-0cm}
\renewcommand{\arraystretch}{1.9}
\begin{tabular}{|p{3.9cm}|p{9.9cm}|}
\hline
\multicolumn{2}{|c|}{\textbf{Caso d’uso: ModificaUtente}} \\ \hline
\textbf{ID} & UC-3 \\ \hline
\textbf{Breve descrizione} & L'Utente visualizza e aggiorna i propri dati personali o la foto profilo. \\ \hline
\textbf{Attori primari} & Utente autenticato \\ \hline
\textbf{Attori secondari} & Database (Cloud Firestore), Cloud Storage (Supabase Storage) \\ \hline
\textbf{Precondizioni} & \begin{enumerate}[leftmargin=14pt,label=\arabic*.,labelsep=0.5em,topsep=0pt,partopsep=0pt,parsep=0pt,itemsep=0pt]
    \item L'Utente è autenticato e si trova nella sezione di gestione del proprio profilo.
\end{enumerate} \\ \hline
\textbf{Sequenza degli eventi principale} &
\begin{enumerate}[leftmargin=14pt,label=\arabic*.,labelsep=0.5em,topsep=0pt,partopsep=0pt,parsep=0pt,itemsep=0pt]
   \item Il sistema mostra i dati attuali precompilati nei rispettivi campi.
   \item L'Utente modifica le informazioni desiderate.
   \item Il sistema valida i nuovi dati.
   \item Il sistema aggiorna le informazioni nei database.
\end{enumerate}\\ \hline
\textbf{Postcondizioni} &
\begin{enumerate}[leftmargin=14pt,label=\arabic*.,labelsep=0.5em,topsep=0pt,partopsep=0pt,parsep=0pt,itemsep=0pt]
\item I dati dell'account Utente sono aggiornati correttamente nel sistema.
\end{enumerate} \\ \hline
\textbf{Sequenza degli eventi alternativa} & Nessuna \\ \hline
\end{tabular}
\end{table}

\newpage

% UC-4: Logout
\begin{table}[H]
\vspace*{-0cm}
\renewcommand{\arraystretch}{1.9}
\begin{tabular}{|p{3.9cm}|p{9.9cm}|}
\hline
\multicolumn{2}{|c|}{\textbf{Caso d’uso: Logout}} \\ \hline
\textbf{ID} & UC-4 \\ \hline
\textbf{Breve descrizione} & L'Utente termina la propria sessione attiva in modo sicuro. \\ \hline
\textbf{Attori primari} & Utente autenticato \\ \hline
\textbf{Attori secondari} & Sistema di Autenticazione (Firebase Auth) \\ \hline
\textbf{Precondizioni} & \begin{enumerate}[leftmargin=14pt,label=\arabic*.,labelsep=0.5em,topsep=0pt,partopsep=0pt,parsep=0pt,itemsep=0pt]
    \item L'Utente è attualmente autenticato nell'applicazione.
\end{enumerate} \\ \hline
\textbf{Sequenza degli eventi principale} &
\begin{enumerate}[leftmargin=14pt,label=\arabic*.,labelsep=0.5em,topsep=0pt,partopsep=0pt,parsep=0pt,itemsep=0pt]
   \item L'Utente naviga nelle impostazioni del profilo e seleziona l'opzione di logout.
   \item Il sistema revoca i token di accesso e distrugge la sessione attiva.
   \item L'Utente viene reindirizzato alla schermata iniziale.
\end{enumerate}\\ \hline
\textbf{Postcondizioni} &
\begin{enumerate}[leftmargin=14pt,label=\arabic*.,labelsep=0.5em,topsep=0pt,partopsep=0pt,parsep=0pt,itemsep=0pt]
\item La sessione è terminata e non è più possibile accedere ai dati protetti senza effettuare un nuovo accesso.
\end{enumerate} \\ \hline
\textbf{Sequenza degli eventi alternativa} & Nessuna \\ \hline
\end{tabular}
\end{table}

% UC-5: Elimina Utente
\begin{table}[H]
\vspace*{-0cm}
\renewcommand{\arraystretch}{1.9}
\begin{tabular}{|p{3.9cm}|p{9.9cm}|}
\hline
\multicolumn{2}{|c|}{\textbf{Caso d’uso: EliminaUtente}} \\ \hline
\textbf{ID} & UC-5 \\ \hline
\textbf{Breve descrizione} & L'Utente richiede l'eliminazione definitiva del proprio account e di tutti i dati associati. \\ \hline
\textbf{Attori primari} & Utente autenticato \\ \hline
\textbf{Attori secondari} & Firebase Auth, Database (Cloud Firestore), Cloud Storage (Supabase Storage) \\ \hline
\textbf{Precondizioni} & \begin{enumerate}[leftmargin=14pt,label=\arabic*.,labelsep=0.5em,topsep=0pt,partopsep=0pt,parsep=0pt,itemsep=0pt]
    \item L'Utente è autenticato e si trova nella sezione di gestione account.
\end{enumerate} \\ \hline
\textbf{Sequenza degli eventi principale} &
\begin{enumerate}[leftmargin=14pt,label=\arabic*.,labelsep=0.5em,topsep=0pt,partopsep=0pt,parsep=0pt,itemsep=0pt]
   \item L'Utente seleziona l'opzione per eliminare l'account.
   \item Il sistema mostra un avviso indicando che l'operazione è irreversibile e richiede un'ulteriore conferma di sicurezza.
   \item L'Utente fornisce la conferma richiesta.
   \item Il sistema procede all’eliminazione dei dati, cancellando le immagini collegate su Supabase Storage, rimuovendo i post, i preferiti e il profilo da Cloud Firestore, ed eliminando l’account da Firebase Auth.
   \item L'Utente viene disconnesso e reindirizzato alla schermata iniziale mostrando un messaggio di conferma.
\end{enumerate}\\ \hline
\textbf{Postcondizioni} &
\begin{enumerate}[leftmargin=14pt,label=\arabic*.,labelsep=0.5em,topsep=0pt,partopsep=0pt,parsep=0pt,itemsep=0pt]
\item L'account e tutti i dati ad esso correlati sono rimossi dal sistema.
\end{enumerate} \\ \hline
\textbf{Sequenza degli eventi alternativa} & \begin{enumerate}[leftmargin=14pt,label=\arabic*.,labelsep=0.5em,topsep=0pt,partopsep=0pt,parsep=0pt,itemsep=0pt]
   \item \texttt{If} l'Utente annulla l'operazione nella finestra di conferma, allora
   \begin{enumerate}[label=\arabic{enumi}.\arabic*.,leftmargin=22pt,labelsep=0.5em,topsep=0pt,partopsep=0pt,parsep=0pt,itemsep=0pt]
       \item Il sistema interrompe la procedura di eliminazione.
       \item L'Utente viene riportato alle impostazioni del profilo senza modifiche ai dati.
   \end{enumerate}
\end{enumerate}\\ \hline
\end{tabular}
\end{table}

% UC-6: Visualizza Utente
\begin{table}[H]
\vspace*{-0cm}
\renewcommand{\arraystretch}{1.9}
\begin{tabular}{|p{3.9cm}|p{9.9cm}|}
\hline
\multicolumn{2}{|c|}{\textbf{Caso d’uso: VisualizzaUtente}} \\ \hline
\textbf{ID} & UC-6 \\ \hline
\textbf{Breve descrizione} & L'Utente consulta la pagina di profilo di un altro utente visualizzarne le informazioni pubbliche e i post pubblicati. \\ \hline
\textbf{Attori primari} & Utente registrato \\ \hline
\textbf{Attori secondari} & Firebase Firestore, Supabase \\ \hline
\textbf{Precondizioni} & \begin{enumerate}[leftmargin=14pt,label=\arabic*.,labelsep=0.5em,topsep=0pt,partopsep=0pt,parsep=0pt,itemsep=0pt]
    \item L'Utente si trova nella scheda di dettaglio di un cucciolo.
\end{enumerate} \\ \hline
\textbf{Sequenza degli eventi principale} &
\begin{enumerate}[leftmargin=14pt,label=\arabic*.,labelsep=0.5em,topsep=0pt,partopsep=0pt,parsep=0pt,itemsep=0pt]
   \item L'Utente seleziona il profilo del proèrietario presente nella scheda del cucciolo.
   \item Il sistema richiede al database i dati del profilo e l'elenco dei cuccioli pubblicati da quell'utente.
   \item Il sistema recupera con successo le informazioni.
   \item Il sistema mostra a schermo la pagina del profilo con i relativi dettagli e post.
\end{enumerate}\\ \hline
\textbf{Postcondizioni} &
\begin{enumerate}[leftmargin=14pt,label=\arabic*.,labelsep=0.5em,topsep=0pt,partopsep=0pt,parsep=0pt,itemsep=0pt]
\item L'Utente visualizza a schermo le informazioni e i post dell'Utente selezionato.
\end{enumerate} \\ \hline
\textbf{Sequenza degli eventi alternativa} & \begin{enumerate}[leftmargin=14pt,label=\arabic*.,labelsep=0.5em,topsep=0pt,partopsep=0pt,parsep=0pt,itemsep=0pt]
   \item \texttt{If} l'Utente cerca il profilo nella barra apposita della schermata principale, allora
   \begin{enumerate}[label=\arabic{enumi}.\arabic*.,leftmargin=22pt,labelsep=0.5em,topsep=0pt,partopsep=0pt,parsep=0pt,itemsep=0pt]
       \item L'Utente seleziona il profilo desiderato dall'elenco dei risultati.
       \item Il sistema richiede i dati al database e mostra a schermo la pagina del profilo.
   \end{enumerate}
\end{enumerate}\\ \hline
\end{tabular}
\end{table}

\clearpage
\newsubsection{Area: Gestione Cucciolo}
Quest'area raggruppa tutti i casi d'uso relativi alla gestione dei post di adozione. Definisce i processi tramite i quali gli utenti possono pubblicare nuovi post relativi ai cuccioli caricandone foto e dettagli,  modificare o rimuovere le proprie inserzioni e, infine, di consultare le schede informative degli animali disponibili.
\vspace{1cm}
\begin{figure}[ht!]
    \centering
    \includegraphics[width=0.9\textwidth]{GestioneCucciolo.png}
    \caption{Gestione Cucciolo}
\end{figure}

% UC-7: Creazione Cucciolo
\begin{table}[H]
\vspace*{-0cm}
\renewcommand{\arraystretch}{1.9}
\begin{tabular}{|p{3.9cm}|p{9.9cm}|}
\hline
\multicolumn{2}{|c|}{\textbf{Caso d’uso: CreaCucciolo}} \\ \hline
\textbf{ID} & UC-7 \\ \hline
\textbf{Breve descrizione} & Un Utente registrato crea e pubblica un nuovo post relativo a un cucciolo fornendo le informazioni richieste e caricando l'immagine. \\ \hline
\textbf{Attori primari} & Utente registrato \\ \hline
\textbf{Attori secondari} & Cloud Storage (Supabase Storage) \\ \hline
\textbf{Precondizioni} & \begin{enumerate}[leftmargin=14pt,label=\arabic*.,labelsep=0.5em,topsep=0pt,partopsep=0pt,parsep=0pt,itemsep=0pt]
   \item L'Utente è autenticato e si trova nella schermata di creazione di un nuovo post.
\end{enumerate} \\ \hline
\textbf{Sequenza degli eventi principale} &
\begin{enumerate}[leftmargin=14pt,label=\arabic*.,labelsep=0.5em,topsep=0pt,partopsep=0pt,parsep=0pt,itemsep=0pt]
  \item L'Utente inserisce le informazioni del cucciolo e seleziona l’immagine.
  \item L'Utente preme il pulsante di pubblicazione.
  \item Il sistema salva il post e i file multimediali nel database.
  \item Il sistema mostra un messaggio di conferma dell'avvenuta creazione e reindirizza l'Utente alla lista cuccioli.
\end{enumerate}\\ \hline
\textbf{Postcondizioni} &
\begin{enumerate}[leftmargin=14pt,label=\arabic*.,labelsep=0.5em,topsep=0pt,partopsep=0pt,parsep=0pt,itemsep=0pt]
\item Il nuovo post del cucciolo è salvato nel sistema ed è visibile all'interno dell'applicazione.
\end{enumerate} \\ \hline
\textbf{Sequenza degli eventi alternativa} & Nessuna \\ \hline
\end{tabular}
\end{table}

% UC-8: Modifica Cucciolo
\begin{table}[H]
\vspace*{-0cm}
\renewcommand{\arraystretch}{1.9}
\begin{tabular}{|p{3.9cm}|p{9.9cm}|}
\hline
\multicolumn{2}{|c|}{\textbf{Caso d’uso: ModificaCucciolo}} \\ \hline
\textbf{ID} & UC-8 \\ \hline
\textbf{Breve descrizione} & Un Utente registrato aggiorna le informazioni o l’immagine di un post precedentemente inserito. \\ \hline
\textbf{Attori primari} & Utente registrato (proprietario del post) \\ \hline
\textbf{Attori secondari} & Cloud Storage (Supabase Storage)
 \\ \hline
\textbf{Precondizioni} & \begin{enumerate}[leftmargin=14pt,label=\arabic*.,labelsep=0.5em,topsep=0pt,partopsep=0pt,parsep=0pt,itemsep=0pt]
   \item L'Utente è autenticato, è il proprietario del post e si trova nella schermata di modifica del cucciolo selezionato.
\end{enumerate} \\ \hline
\textbf{Sequenza degli eventi principale} &
\begin{enumerate}[leftmargin=14pt,label=\arabic*.,labelsep=0.5em,topsep=0pt,partopsep=0pt,parsep=0pt,itemsep=0pt]
  \item Il sistema precompila il modulo con i dati e i file attualmente presenti.
  \item L'Utente modifica uno o più campi.
  \item L'Utente preme il pulsante di salvataggio delle modifiche.
  \item Il sistema verifica la correttezza dei dati aggiornati.
  \item Il sistema aggiorna le informazioni e le immagini sul database.
  \item Il sistema mostra un messaggio di successo e reindirizza l'Utente alla lista cuccioli.
\end{enumerate}\\ \hline
\textbf{Postcondizioni} &
\begin{enumerate}[leftmargin=14pt,label=\arabic*.,labelsep=0.5em,topsep=0pt,partopsep=0pt,parsep=0pt,itemsep=0pt]
\item Le informazioni del post del cucciolo risultano aggiornate con successo nel database.
\end{enumerate} \\ \hline
\textbf{Sequenza degli eventi alternativa} & Nessuna \\ \hline
\end{tabular}
\end{table}

% UC-9: Eliminazione Cucciolo
\begin{table}[H]
\vspace*{-0cm}
\renewcommand{\arraystretch}{1.9}
\begin{tabular}{|p{3.9cm}|p{9.9cm}|}
\hline
\multicolumn{2}{|c|}{\textbf{Caso d’uso: EliminaCucciolo}} \\ \hline
\textbf{ID} & UC-9 \\ \hline
\textbf{Breve descrizione} & Un Utente registrato rimuove in modo permanente un post di un cucciolo di sua proprietà dal sistema. \\ \hline
\textbf{Attori primari} & Utente registrato (proprietario del post) \\ \hline
\textbf{Attori secondari} & Cloud Storage (Supabase Storage)
 \\ \hline
\textbf{Precondizioni} & \begin{enumerate}[leftmargin=14pt,label=\arabic*.,labelsep=0.5em,topsep=0pt,partopsep=0pt,parsep=0pt,itemsep=0pt]
   \item L'Utente è autenticato, è il proprietario del post, e si trova nella lista dei propri post.
\end{enumerate} \\ \hline
\textbf{Sequenza degli eventi principale} &
\begin{enumerate}[leftmargin=14pt,label=\arabic*.,labelsep=0.5em,topsep=0pt,partopsep=0pt,parsep=0pt,itemsep=0pt]
  \item L'Utente seleziona l'opzione per eliminare il post del cucciolo.
  \item Il sistema richiede una conferma all'Utente prima di procedere.
  \item L'Utente conferma l'eliminazione.
  \item Il sistema rimuove definitivamente il post e le relative immagini dal database.
  \item Il sistema mostra un messaggio di conferma dell'eliminazione.
\end{enumerate}\\ \hline
\textbf{Postcondizioni} &
\begin{enumerate}[leftmargin=14pt,label=\arabic*.,labelsep=0.5em,topsep=0pt,partopsep=0pt,parsep=0pt,itemsep=0pt]
\item Il post del cucciolo è rimosso dal database e non è più visibile o ricercabile nell'applicazione.
\end{enumerate} \\ \hline
\textbf{Sequenza degli eventi alternativa} & \begin{enumerate}[leftmargin=14pt,label=\arabic*.,labelsep=0.5em,topsep=0pt,partopsep=0pt,parsep=0pt,itemsep=0pt]
  \item \texttt{If} l'Utente annulla la richiesta nella finestra di conferma, allora
  \begin{enumerate}[label=\arabic{enumi}.\arabic*.,leftmargin=22pt,labelsep=0.5em,topsep=0pt,partopsep=0pt,parsep=0pt,itemsep=0pt]
      \item Il sistema chiude il dialogo di conferma senza cancellare i dati.
      \item Il post del cucciolo rimane inalterato nel sistema.
  \end{enumerate}
\end{enumerate}\\ \hline
\end{tabular}
\end{table}

% UC-10: Visualizza Cucciolo
\begin{table}[H]
\vspace*{-0cm}
\renewcommand{\arraystretch}{1.9}
\begin{tabular}{|p{3.9cm}|p{9.9cm}|}
\hline
\multicolumn{2}{|c|}{\textbf{Caso d’uso: VisualizzaCucciolo}} \\ \hline
\textbf{ID} & UC-10 \\ \hline
\textbf{Breve descrizione} & Un Utente consulta la scheda di dettaglio di un cucciolo per visualizzarne le informazioni e le foto. \\ \hline
\textbf{Attori primari} & Utente registrato\\ \hline
\textbf{Attori secondari} & Supabase Storage \\ \hline
\textbf{Precondizioni} & \begin{enumerate}[leftmargin=14pt,label=\arabic*.,labelsep=0.5em,topsep=0pt,partopsep=0pt,parsep=0pt,itemsep=0pt]
   \item L'Utente si trova nella schermata principale o nel profilo di un utente.
\end{enumerate} \\ \hline
\textbf{Sequenza degli eventi principale} &
\begin{enumerate}[leftmargin=14pt,label=\arabic*.,labelsep=0.5em,topsep=0pt,partopsep=0pt,parsep=0pt,itemsep=0pt]
  \item L'Utente preme sulle informazioni del cucciolo.
  \item Il sistema richiede al database i dati dettagliati e i contenuti multimediali del cucciolo.
  \item Il sistema recupera con successo le informazioni.
  \item Il sistema mostra la schermata con i dettagli del cucciolo.
\end{enumerate}\\ \hline
\textbf{Postcondizioni} &
\begin{enumerate}[leftmargin=14pt,label=\arabic*.,labelsep=0.5em,topsep=0pt,partopsep=0pt,parsep=0pt,itemsep=0pt]
\item L'Utente visualizza a schermo le informazioni dettagliate del cucciolo richiesto.
\end{enumerate} \\ \hline
\textbf{Sequenza degli eventi alternativa} & \begin{enumerate}[leftmargin=14pt,label=\arabic*.,labelsep=0.5em,topsep=0pt,partopsep=0pt,parsep=0pt,itemsep=0pt]
   \item \texttt{If} l'Utente si trova nella lista dei preferiti, allora
   \begin{enumerate}[label=\arabic{enumi}.\arabic*.,leftmargin=22pt,labelsep=0.5em,topsep=0pt,partopsep=0pt,parsep=0pt,itemsep=0pt]
       \item L'Utente preme sull'anteprima del cucciolo per accedere ai dettagli.
       \item Il sistema richiede al database i dati
       \item Il sistema mostra la schermata con i dettagli del cucciolo.
   \end{enumerate}
\end{enumerate}\\ \hline
\end{tabular}
\end{table}


\clearpage
\newsubsection{Area: Gestione Interazioni}
Quest'area raggruppa i casi d'uso relativi all'interazione con i post e con i proprietari dei cuccioli. Definisce i processi tramite i quali gli utenti possono salvare gli animali nella propria lista dei preferiti o avviare rapidamente un chiamata con altri utenti o rifugi, e infine descrive il meccanismo automatico con cui il sistema calcola e aggiorna il numero totale di "Mi piace" ricevuti da ciascun annuncio.
\vspace{1cm}
\begin{figure}[ht!]
    \centering
    \includegraphics[width=1\textwidth]{GestioneInterazioni.png}
    \caption{Gestione Interazioni}
\end{figure}

% UC-11: Aggiungi Preferiti
\begin{table}[H]
\vspace*{-0cm}
\renewcommand{\arraystretch}{1.9}
\begin{tabular}{|p{3.9cm}|p{9.9cm}|}
\hline
\multicolumn{2}{|c|}{\textbf{Caso d’uso: AggiungiPreferiti}} \\ \hline
\textbf{ID} & UC-11 \\ \hline
\textbf{Breve descrizione} & Un Utente registrato aggiunge un cucciolo alla propria lista dei preferiti. \\ \hline
\textbf{Attori primari} & Utente registrato \\ \hline
\textbf{Attori secondari} & Firebase Firestore \\ \hline
\textbf{Precondizioni} & \begin{enumerate}[leftmargin=14pt,label=\arabic*.,labelsep=0.5em,topsep=0pt,partopsep=0pt,parsep=0pt,itemsep=0pt]
    \item L'Utente è autenticato e si trova nella schermata principale.
\end{enumerate} \\ \hline
\textbf{Sequenza degli eventi principale} &
\begin{enumerate}[leftmargin=14pt,label=\arabic*.,labelsep=0.5em,topsep=0pt,partopsep=0pt,parsep=0pt,itemsep=0pt]
   \item L'Utente ggiunge un cucciolo ai preferiti effettuando un gesto di scorrimento verso destra (Swipe to Favorite) sulla scheda del cucciolo.
   \item Il sistema memorizza l'associazione tra l'Utente e il cucciolo nel database.
   \item Il sistema aggiorna l'interfaccia visiva per confermare l'avvenuto salvataggio.
\end{enumerate}\\ \hline
\textbf{Postcondizioni} &
\begin{enumerate}[leftmargin=14pt,label=\arabic*.,labelsep=0.5em,topsep=0pt,partopsep=0pt,parsep=0pt,itemsep=0pt]
\item Il cucciolo risulta salvato nella lista dei preferiti dell'Utente.
\end{enumerate} \\ \hline
\textbf{Sequenza degli eventi alternativa} & \begin{enumerate}[leftmargin=14pt,label=\arabic*.,labelsep=0.5em,topsep=0pt,partopsep=0pt,parsep=0pt,itemsep=0pt]
   \item \texttt{If} l'operazione viene eseguita dalla pagina di profilo di un altro Utente, allora
   \begin{enumerate}[label=\arabic{enumi}.\arabic*.,leftmargin=22pt,labelsep=0.5em,topsep=0pt,partopsep=0pt,parsep=0pt,itemsep=0pt]
       \item L'Utente seleziona l'icona dei preferiti (cuore) posta sull'anteprima del cucciolo.
       \item Il sistema registra l'aggiunta ai preferiti nel database.
       \item Il sistema aggiorna l'interfaccia attivando visivamente l'icona.
   \end{enumerate}
\end{enumerate}\\ \hline
\end{tabular}
\end{table}

% UC-12: Rimozione Preferiti
\begin{table}[H]
\vspace*{-0cm}
\renewcommand{\arraystretch}{1.9}
\begin{tabular}{|p{3.9cm}|p{9.9cm}|}
\hline
\multicolumn{2}{|c|}{\textbf{Caso d’uso: RimuoviPreferiti}} \\ \hline
\textbf{ID} & UC-12 \\ \hline
\textbf{Breve descrizione} & Un Utente registrato rimuove un cucciolo precedentemente salvato dalla propria lista dei preferiti. \\ \hline
\textbf{Attori primari} & Utente registrato \\ \hline
\textbf{Attori secondari} & Firebase Firestore \\ \hline
\textbf{Precondizioni} & \begin{enumerate}[leftmargin=14pt,label=\arabic*.,labelsep=0.5em,topsep=0pt,partopsep=0pt,parsep=0pt,itemsep=0pt]
    \item L'Utente è autenticato e si trova nella sezione dei preferiti.
\end{enumerate} \\ \hline
\textbf{Sequenza degli eventi principale} &
\begin{enumerate}[leftmargin=14pt,label=\arabic*.,labelsep=0.5em,topsep=0pt,partopsep=0pt,parsep=0pt,itemsep=0pt]
   \item L'Utente rimuove un cucciolo dai preferiti selezionando l'icona del cuore sull'anteprima del cucciolo.
   \item Il sistema rimuove l'associazione tra l'Utente e il cucciolo nel database.
   \item Il sistema aggiorna l'interfaccia visiva per confermare l'avvenuta rimozione.
\end{enumerate}\\ \hline
\textbf{Postcondizioni} &
\begin{enumerate}[leftmargin=14pt,label=\arabic*.,labelsep=0.5em,topsep=0pt,partopsep=0pt,parsep=0pt,itemsep=0pt]
\item Il cucciolo è rimosso definitivamente dalla lista dei preferiti dell'Utente.
\end{enumerate} \\ \hline
\textbf{Sequenza degli eventi alternativa} & \begin{enumerate}[leftmargin=14pt,label=\arabic*.,labelsep=0.5em,topsep=0pt,partopsep=0pt,parsep=0pt,itemsep=0pt]
   \item \texttt{If} l'operazione viene eseguita dalla pagina di profilo di un altro Utente, allora
   \begin{enumerate}[label=\arabic{enumi}.\arabic*.,leftmargin=22pt,labelsep=0.5em,topsep=0pt,partopsep=0pt,parsep=0pt,itemsep=0pt]
       \item L'Utente seleziona l'icona dei preferiti posta sull'anteprima del cucciolo.
       \item Il sistema registra la rimozione dai preferiti nel database.
       \item Il sistema aggiorna l'interfaccia disattivando visivamente l'icona.
   \end{enumerate}
\end{enumerate}\\ \hline
\end{tabular}
\end{table}

% UC-13: Calcolo Mi Piace
\begin{table}[H]
\vspace*{-0cm}
\renewcommand{\arraystretch}{1.9}
\begin{tabular}{|p{3.9cm}|p{9.9cm}|}
\hline
\multicolumn{2}{|c|}{\textbf{Caso d’uso: CalcoloMiPiace}} \\ \hline
\textbf{ID} & UC-13 \\ \hline
\textbf{Breve descrizione} & Il sistema calcola e aggiorna il numero totale di "Mi piace" ricevuti dal post di un cucciolo. \\ \hline
\textbf{Attori primari} & Utente registrato \\ \hline
\textbf{Attori secondari} & Firebase Firestore \\ \hline
\textbf{Precondizioni} & \begin{enumerate}[leftmargin=14pt,label=\arabic*.,labelsep=0.5em,topsep=0pt,partopsep=0pt,parsep=0pt,itemsep=0pt]
    \item L'Utente sta visualizzando il post di un cucciolo nella schermata principale.
\end{enumerate} \\ \hline
\textbf{Sequenza degli eventi principale} &
\begin{enumerate}[leftmargin=14pt,label=\arabic*.,labelsep=0.5em,topsep=0pt,partopsep=0pt,parsep=0pt,itemsep=0pt]
   \item Il sistema richiede al database il totale delle interazioni registrate per il post.
   \item Il database restituisce il conteggio aggiornato.
   \item Il sistema calcola la somma totale e la formatta per la visualizzazione.
   \item Il sistema mostra a schermo il numero totale di "Mi piace" accanto all'apposita icona.
\end{enumerate}\\ \hline
\textbf{Postcondizioni} &
\begin{enumerate}[leftmargin=14pt,label=\arabic*.,labelsep=0.5em,topsep=0pt,partopsep=0pt,parsep=0pt,itemsep=0pt]
\item Il totale dei "Mi piace" viene correttamente calcolato e mostrato nell'interfaccia del post.
\end{enumerate} \\ \hline
\textbf{Sequenza degli eventi alternativa} & Nessuna \\ \hline
\end{tabular}
\end{table}

% UC-14: Chiamata Telefonica
\begin{table}[H]
\vspace*{-0cm}
\renewcommand{\arraystretch}{1.9}
\begin{tabular}{|p{3.9cm}|p{9.9cm}|}
\hline
\multicolumn{2}{|c|}{\textbf{Caso d’uso: ChiamataTelefonica}} \\ \hline
\textbf{ID} & UC-14 \\ \hline
\textbf{Breve descrizione} & L'Utente avvia direttamente una chiamata telefonica verso il recapito del proprietario del cucciolo. \\ \hline
\textbf{Attori primari} & Utente \\ \hline
\textbf{Attori secondari} & Firebase Firestore \\ \hline
\textbf{Precondizioni} & \begin{enumerate}[leftmargin=14pt,label=\arabic*.,labelsep=0.5em,topsep=0pt,partopsep=0pt,parsep=0pt,itemsep=0pt]
    \item L'Utente visualizza la scheda di dettaglio di un cucciolo in cui è presente il recapito di contatto.
\end{enumerate} \\ \hline
\textbf{Sequenza degli eventi principale} &
\begin{enumerate}[leftmargin=14pt,label=\arabic*.,labelsep=0.5em,topsep=0pt,partopsep=0pt,parsep=0pt,itemsep=0pt]
   \item L'Utente preme il numero di telefono nella scheda del cucciolo.
   \item Il sistema recupera il numero di telefono del proprietario associato al post.
   \item Il sistema richiama le funzionalità del sistema operativo dello smartphone 
   \item Il sistema chiede conferma per l'avvio della chiamata
\end{enumerate}\\ \hline
\textbf{Postcondizioni} &
\begin{enumerate}[leftmargin=14pt,label=\arabic*.,labelsep=0.5em,topsep=0pt,partopsep=0pt,parsep=0pt,itemsep=0pt]
   \item L'Utente avvia la chiamata verso il proprietario.
\end{enumerate} \\ \hline
\textbf{Sequenza degli eventi alternativa} & Nessuna \\ \hline
\end{tabular}
\end{table}

\clearpage
\newsubsection{Area: Gestione Navigazione}
Quest'area raggruppa le funzionalità tutti i casi d'uso relativi all'esplorazione del database dei post. Definisce i processi tramite i quali gli utenti possonoeffettuare ricerche testuali mirate tramite parole chiave e di applicare filtri avanzati per restringere i risultati mostrati a schermo, facilitando così l'individuazione del cucciolo più in linea con le proprie preferenze. 
\vspace{1cm}
\begin{figure}[ht!]
    \centering
    \includegraphics[width=1\textwidth]{GestioneNavigazione.png}
    \caption{Gestione Navigazione}
\end{figure}

% UC-15: Cerca
\begin{table}[H]
\vspace*{-0cm}
\renewcommand{\arraystretch}{1.9}
\begin{tabular}{|p{3.9cm}|p{9.9cm}|}
\hline
\multicolumn{2}{|c|}{\textbf{Caso d’uso: Cerca}} \\ \hline
\textbf{ID} & UC-15 \\ \hline
\textbf{Breve descrizione} & L'Utente inserisce una o più parole chiave per cercare utenti o cuccioli nell'applicazione. \\ \hline
\textbf{Attori primari} & Utente (registrato o non registrato) \\ \hline
\textbf{Attori secondari} & Database Remoto (es. Firebase Firestore / Supabase) \\ \hline
\textbf{Precondizioni} & \begin{enumerate}[leftmargin=14pt,label=\arabic*.,labelsep=0.5em,topsep=0pt,partopsep=0pt,parsep=0pt,itemsep=0pt]
    \item L'Utente si trova nella schermata principale.
\end{enumerate} \\ \hline
\textbf{Sequenza degli eventi principale} &
\begin{enumerate}[leftmargin=14pt,label=\arabic*.,labelsep=0.5em,topsep=0pt,partopsep=0pt,parsep=0pt,itemsep=0pt]
   \item L'Utente digita la parola chiave nella barra di ricerca.
   \item L'Utente avvia la ricerca.
   \item Il sistema richiede al database i dati che corrispondono al testo inserito.
   \item Il sistema recupera i dati e mostra a schermo l'elenco dei risultati trovati.
\end{enumerate}\\ \hline
\textbf{Postcondizioni} &
\begin{enumerate}[leftmargin=14pt,label=\arabic*.,labelsep=0.5em,topsep=0pt,partopsep=0pt,parsep=0pt,itemsep=0pt]
\item L'applicazione mostra la lista degli utenti o dei cuccioli corrispondenti al testo cercato.
\end{enumerate} \\ \hline
\textbf{Sequenza degli eventi alternativa} & \begin{enumerate}[leftmargin=14pt,label=\arabic*.,labelsep=0.5em,topsep=0pt,partopsep=0pt,parsep=0pt,itemsep=0pt]
   \item \texttt{If} l'Utente si trova nella schermata dei propri post o nella lista dei preferiti, allora
   \begin{enumerate}[label=\arabic{enumi}.\arabic*.,leftmargin=22pt,labelsep=0.5em,topsep=0pt,partopsep=0pt,parsep=0pt,itemsep=0pt]
       \item L'Utente digita la parola chiave nell'apposita barra di ricerca locale.
       \item Il sistema filtra l'elenco specifico richiedendo al database i risultati pertinenti.
       \item Il sistema mostra a schermo l'elenco dei cuccioli filtrati.
   \end{enumerate}
\end{enumerate}\\ \hline
\end{tabular}
\end{table}


% UC-16: Filtra
\begin{table}[H]
\vspace*{-0cm}
\renewcommand{\arraystretch}{1.9}
\begin{tabular}{|p{3.9cm}|p{9.9cm}|}
\hline
\multicolumn{2}{|c|}{\textbf{Caso d’uso: Filtra}} \\ \hline
\textbf{ID} & UC-16 \\ \hline
\textbf{Breve descrizione} & L'Utente seleziona determinati parametri per restringere la lista degli utenti o cuccioli da visualizzare. \\ \hline
\textbf{Attori primari} & Utente (registrato o non registrato) \\ \hline
\textbf{Attori secondari} & Firebase Firestore, Supabase \\ \hline
\textbf{Precondizioni} & \begin{enumerate}[leftmargin=14pt,label=\arabic*.,labelsep=0.5em,topsep=0pt,partopsep=0pt,parsep=0pt,itemsep=0pt]
    \item L'Utente è nella schermata principale ed apre il menu dei filtri.
\end{enumerate} \\ \hline
\textbf{Sequenza degli eventi principale} &
\begin{enumerate}[leftmargin=14pt,label=\arabic*.,labelsep=0.5em,topsep=0pt,partopsep=0pt,parsep=0pt,itemsep=0pt]
   \item L'Utente seleziona uno o più criteri di filtraggio.
   \item L'Utente preme il pulsante per applicare i filtri.
   \item Il sistema richiede al database unicamente i post che soddisfano tutti i parametri selezionati.
   \item Il sistema aggiorna l'interfaccia mostrata all'Utente con la nuova lista filtrata.
\end{enumerate}\\ \hline
\textbf{Postcondizioni} &
\begin{enumerate}[leftmargin=14pt,label=\arabic*.,labelsep=0.5em,topsep=0pt,partopsep=0pt,parsep=0pt,itemsep=0pt]
\item L'elenco visualizzato contiene esclusivamente gli utenti o i cuccioli che rispettano i filtri applicati.
\end{enumerate} \\ \hline
\textbf{Sequenza degli eventi alternativa} & \begin{enumerate}[leftmargin=14pt,label=\arabic*.,labelsep=0.5em,topsep=0pt,partopsep=0pt,parsep=0pt,itemsep=0pt]
   \item \texttt{If} l'Utente si trova nella sezione dei propri post o nella lista dei preferiti, allora
   \begin{enumerate}[label=\arabic{enumi}.\arabic*.,leftmargin=22pt,labelsep=0.5em,topsep=0pt,partopsep=0pt,parsep=0pt,itemsep=0pt]
       \item L'Utente apre il menu dei filtri e seleziona i criteri desiderati.
       \item Il sistema filtra l'elenco specifico (post personali o preferiti) richiedendo al database solo i risultati pertinenti.
       \item Il sistema mostra a schermo l'elenco filtrato.
   \end{enumerate}
\end{enumerate}\\ \hline
\end{tabular}
\end{table}

\clearpage
\newsubsection{Area: Gestione Notifiche}
Quest'area raggruppa i casi d'uso dedicati al sistema di comunicazione e di avvisi in tempo reale. Gestisce la generazione e l'invio automatico delle notifiche push a fronte di eventi specifici di interesse per l'utente, consentendo inoltre la consultazione della cronologia degli avvisi e l'eliminazione di quest'ultimi per mantenere pulito il proprio centro notifiche.
\vspace{1cm}
\begin{figure}[ht!]
    \centering
    \includegraphics[width=1\textwidth]{GestioneNotifiche.png}
    \caption{Gestione Notifiche}
\end{figure}

% UC-17: Genera Notifica
\begin{table}[H]
\vspace*{-0cm}
\renewcommand{\arraystretch}{1.9}
\begin{tabular}{|p{3.9cm}|p{9.9cm}|}
\hline
\multicolumn{2}{|c|}{\textbf{Caso d’uso: GeneraNotifica}} \\ \hline
\textbf{ID} & UC-17 \\ \hline
\textbf{Breve descrizione} & Il sistema genera e invia automaticamente una notifica all'Utente al verificarsi di un evento di suo interesse. \\ \hline
\textbf{Attori primari} & Sistema \\ \hline
\textbf{Attori secondari} & Firebase Cloud Messaging \\ \hline
\textbf{Precondizioni} & \begin{enumerate}[leftmargin=14pt,label=\arabic*.,labelsep=0.5em,topsep=0pt,partopsep=0pt,parsep=0pt,itemsep=0pt]
    \item Si verifica un evento nel sistema che richiede l'invio di una notifica a un Utente destinatario.
\end{enumerate} \\ \hline
\textbf{Sequenza degli eventi principale} &
\begin{enumerate}[leftmargin=14pt,label=\arabic*.,labelsep=0.5em,topsep=0pt,partopsep=0pt,parsep=0pt,itemsep=0pt]
   \item Il sistema rileva l'evento scatenante.
   \item Il sistema compone il titolo, il testo e i dati della notifica.
   \item Il sistema salva la notifica nel database dell'Utente destinatario.
   \item Il sistema invia la notifica push al dispositivo dell'Utente tramite il servizio dedicato.
\end{enumerate}\\ \hline
\textbf{Postcondizioni} &
\begin{enumerate}[leftmargin=14pt,label=\arabic*.,labelsep=0.5em,topsep=0pt,partopsep=0pt,parsep=0pt,itemsep=0pt]
\item La notifica viene registrata nella sezione Notifiche e recapitata al suo dispositivo.
\end{enumerate} \\ \hline
\textbf{Sequenza degli eventi alternativa} & Nessuna \\ \hline
\end{tabular}
\end{table}

% UC-18: Visualizza Notifica
\begin{table}[H]
\vspace*{-0cm}
\renewcommand{\arraystretch}{1.9}
\begin{tabular}{|p{3.9cm}|p{9.9cm}|}
\hline
\multicolumn{2}{|c|}{\textbf{Caso d’uso: VisualizzaNotifica}} \\ \hline
\textbf{ID} & UC-18 \\ \hline
\textbf{Breve descrizione} & L'Utente registrato accede al centro notifiche, consulta l'elenco delle notifiche ricevute. \\ \hline
\textbf{Attori primari} & Utente registrato \\ \hline
\textbf{Attori secondari} & Firebase Firestore \\ \hline
\textbf{Precondizioni} & \begin{enumerate}[leftmargin=14pt,label=\arabic*.,labelsep=0.5em,topsep=0pt,partopsep=0pt,parsep=0pt,itemsep=0pt]
    \item L'Utente è autenticato e seleziona la sezione "Notifiche" dell'applicazione.
\end{enumerate} \\ \hline
\textbf{Sequenza degli eventi principale} &
\begin{enumerate}[leftmargin=14pt,label=\arabic*.,labelsep=0.5em,topsep=0pt,partopsep=0pt,parsep=0pt,itemsep=0pt]
   \item L'Utente entra nella schermata delle notifiche.
   \item Il sistema richiede al database la lista delle notifiche dell'Utente.
   \item Il sistema mostra le notifiche ordinate dalla più recente.
\end{enumerate}\\ \hline
\textbf{Postcondizioni} &
\begin{enumerate}[leftmargin=14pt,label=\arabic*.,labelsep=0.5em,topsep=0pt,partopsep=0pt,parsep=0pt,itemsep=0pt]
\item L'Utente prende visualizza la notifica.
\end{enumerate} \\ \hline
\textbf{Sequenza degli eventi alternativa} & \begin{enumerate}[leftmargin=14pt,label=\arabic*.,labelsep=0.5em,topsep=0pt,partopsep=0pt,parsep=0pt,itemsep=0pt]
   \item \texttt{If} non sono presenti notifiche nell'account dell'Utente, allora
   \begin{enumerate}[label=\arabic{enumi}.\arabic*.,leftmargin=22pt,labelsep=0.5em,topsep=0pt,partopsep=0pt,parsep=0pt,itemsep=0pt]
       \item Il sistema mostra una schermata vuota informativa.
   \end{enumerate}
\end{enumerate}\\ \hline
\end{tabular}
\end{table}

% UC-19: Elimina Notifica
\begin{table}[H]
\vspace*{-0cm}
\renewcommand{\arraystretch}{1.9}
\begin{tabular}{|p{3.9cm}|p{9.9cm}|}
\hline
\multicolumn{2}{|c|}{\textbf{Caso d’uso: EliminaNotifica}} \\ \hline
\textbf{ID} & UC-19 \\ \hline
\textbf{Breve descrizione} & L'Utente registrato le notifiche dal proprio elenco. \\ \hline
\textbf{Attori primari} & Utente registrato \\ \hline
\textbf{Attori secondari} & Firebase Firestone \\ \hline
\textbf{Precondizioni} & \begin{enumerate}[leftmargin=14pt,label=\arabic*.,labelsep=0.5em,topsep=0pt,partopsep=0pt,parsep=0pt,itemsep=0pt]
    \item L'Utente è autenticato e si trova all'interno della schermata "Notifiche".
\end{enumerate} \\ \hline
\textbf{Sequenza degli eventi principale} &
\begin{enumerate}[leftmargin=14pt,label=\arabic*.,labelsep=0.5em,topsep=0pt,partopsep=0pt,parsep=0pt,itemsep=0pt]
   \item L'Utente seleziona l'opzione di eliminazione.
   \item Il sistema rimuove la notifica dal database.
   \item Il sistema aggiorna l'interfaccia rimuovendo la notifica dall'elenco.
\end{enumerate}\\ \hline
\textbf{Postcondizioni} &
\begin{enumerate}[leftmargin=14pt,label=\arabic*.,labelsep=0.5em,topsep=0pt,partopsep=0pt,parsep=0pt,itemsep=0pt]
\item La notifica eliminata viene cancellata permanentemente dal database dell'Utente.
\end{enumerate} \\ \hline
\textbf{Sequenza degli eventi alternativa} & Nessuna \\ \hline
\end{tabular}
\end{table}

\newchapter{Database e Storage}

La gestione dei dati nell'applicazione \textbf{Paws} adotta un'architettura distribuita e multi-database, progettata per garantire elevata reattività, sicurezza e scalabilità. A differenza di un approccio monolitico basato su un unico motore relazionale locale o cloud, Paws integra in modo sinergico diverse tecnologie specializzate di database e storage, sia per l'ambiente nativo Android in Kotlin che per la versione cross-platform in Flutter:

\begin{itemize}
    \item \textbf{Database Cloud NoSQL (Firebase Cloud Firestore):} gestisce la persistenza dei dati di dominio in tempo reale (utenti, annunci di cuccioli, preferiti e notifiche).
    \item \textbf{Cloud Object Storage (Supabase Storage):} gestisce l'archiviazione efficiente dei file multimediali ad alta risoluzione (immagini del profilo e foto degli annunci).
    \item \textbf{Database di Autenticazione (Firebase Auth):} memorizza le credenziali e gestisce l'identità degli utenti generando gli identificatori univoci (\texttt{UID}).
    \item \textbf{Local Storage (SharedPreferences):} memorizza in locale piccoli quantitativi di dati a chiave-valore per lo stato della sessione e le preferenze grafiche.
\end{itemize}

\newsection{Database Remoto NoSQL (Cloud Firestore)}

Il database principale dell'applicazione è \textbf{Firebase Cloud Firestore}, un database documentale NoSQL basato su collezioni e documenti JSON-like. Firestore è stato scelto per la sua capacità di offrire la sincronizzazione in tempo reale tramite \textit{WebSocket} e \textit{Real-Time Listeners}, consentendo all'interfaccia utente di aggiornarsi istantaneamente al variare dei dati sul cloud.

\noindent Il database è organizzato nelle seguenti collezioni principali:

\begin{itemize}
    \item \textbf{\texttt{users}:} contiene i profili informativi degli utenti e delle strutture di adozione (rifugi/canili). Ogni documento ha come ID il campo \texttt{uid} derivato da Firebase Auth e comprende: \texttt{name}, \texttt{surname}, \texttt{email}, \texttt{phone}, \texttt{role} (privato o rifugio), \texttt{address}, \texttt{description} e \texttt{avatarUrl}.
    \item \textbf{\texttt{posts}:} raccoglie gli annunci dei cuccioli in cerca di casa. Ciascun documento memorizza: \texttt{postId}, \texttt{uid} (proprietario dell'annuncio), \texttt{name} del cucciolo, \texttt{type} (specie animale), \texttt{breed} (razza), \texttt{age}, \texttt{sex}, \texttt{caption} (descrizione), \texttt{imageUrl} (link Supabase) e \texttt{timestamp} di creazione.
    \item \textbf{\texttt{favorites}:} rappresenta la relazione molti-a-molti tra gli utenti e i post salvati nei preferiti. Mantiene l'associazione \texttt{userId} -- \texttt{postId}, consentendo il recupero rapido nella schermata dedicata e supportando le gesture native di \textit{Swipe-to-Favorite}.
    \item \textbf{\texttt{notifications}:} traccia le notifiche di sistema ed interazione inviate agli utenti. Registra l'ID del destinatario (\texttt{recipientId}), l'ID del mittente (\texttt{senderId}), il post di riferimento, il tipo di notifica, lo stato di lettura (\texttt{isRead}) e il marcatore temporale.
\end{itemize}

\newsection{Cloud Object Storage (Supabase Storage)}

Per evitare l'appesantimento del database NoSQL con stringhe Base64 o oggetti binari di grandi dimensioni (BLOB), l'architettura di Paws disaccoppia la gestione dei dati testuali da quella dei contenuti multimediali affidandosi a \textbf{Supabase Storage}.

\noindent Supabase mette a disposizione bucket di storage dedicati:
\begin{itemize}
    \item \textbf{Bucket \texttt{profiles}:} dedicato alle immagini di profilo degli utenti e dei canili.
    \item \textbf{Bucket \texttt{puppies}:} dedicato al caricamento delle foto dei cuccioli associate agli annunci.
\end{itemize}

\noindent Quando un utente carica un'immagine (tramite la \texttt{CreatePostFragment} in Kotlin o l'equivalente in Flutter), il file binario viene inviato tramite le API REST/SDK di Supabase Storage. Supabase restituisce un URL pubblico sicuro, il quale viene poi salvato nel relativo documento Firestore (\texttt{imageUrl} o \texttt{avatarUrl}). Lato client, le librerie di visualizzazione quali \textbf{Glide} (in Android Kotlin) e \textbf{\texttt{cached\_network\_image}} (in Flutter) caricano le immagini in modo asincrono sfruttando tecniche di caching in memoria e su disco.

\newsection{Local Storage (SharedPreferences)}

Accanto ai database di rete cloud, l'applicazione sfrutta il meccanismo di memorizzazione locale a chiave-valore offerto da \textbf{SharedPreferences} (in Android Native) e dal pacchetto \textbf{\texttt{shared\_preferences}} (in Flutter).

\noindent Questo storage locale viene impiegato principalmente per:
\begin{itemize}
    \item Salvare lo stato della sessione e l'ID utente corrente (\texttt{UID}) dopo il primo accesso.
    \item Memorizzare le preferenze grafiche dell'interfaccia utente (come la scelta del tema chiaro o scuro).
    \item Salvare temporaneamente il flag di completamento dell'onboarding.
\end{itemize}

\noindent In questo modo, all'avvio dell'applicazione è possibile determinare istantaneamente se reindirizzare l'utente alla schermata di Login o caricare la Home, evitando schermate di caricamento vuote e senza attendere l'inizializzazione dei database cloud o delle chiamate di rete.

\newsection{Sincronizzazione Offline e Persistenza Locale}

Una delle caratteristiche fondamentali dell'architettura multi-database di Paws è la capacità di operare anche in assenza temporanea di connettività internet. A differenza delle architetture tradizionali basate su SQL che richiedono la scrittura manuale di tabelle di sync ed attributi come \texttt{isSynced}, \textbf{Firebase Cloud Firestore} integra nativamente un motore di persistenza locale basato su cache interna.

\noindent Il funzionamento del meccanismo di sincronizzazione offline segue questo comportamento:
\begin{itemize}
    \item \textbf{Scrittura offline:} quando l'utente crea un annuncio o aggiunge un cucciolo ai preferiti senza rete internet, lo SDK di Firestore memorizza l'operazione nel database SQLite/LevelDB locale del dispositivo e aggiorna immediatamente la UI (\textit{Optimistic UI Updates}).
    \item \textbf{Ripristino connettività:} al rilevamento del segnale di rete, lo SDK invia in automatico le modifiche in coda ai server cloud Firebase, garantendo l'allineamento definitivo dei dati in background.
\end{itemize}

\newsection{Politiche di Sicurezza (Firestore Security Rules)}

Per proteggere l'integrità dei dati salvati nei database cloud e prevenire accessi o modifiche non autorizzate, sono state configurate apposite regole di sicurezza dichiarative su Cloud Firestore (\textit{Security Rules}):

\begin{scriptsize}
\estiloPython
\begin{lstlisting}[caption={Regole di Sicurezza Cloud Firestore}, label=lst:firestore_rules]
rules_version = '2';
service cloud.firestore {
  match /databases/{database}/documents {
    
    // Regole per la collezione utenti
    match /users/{userId} {
      allow read: if request.auth != null;
      allow create, update: if request.auth != null && request.auth.uid == userId;
      allow delete: if request.auth != null && request.auth.uid == userId;
    }
    
    // Regole per la collezione post di adozione cuccioli
    match /posts/{postId} {
      allow read: if request.auth != null;
      allow create: if request.auth != null && request.resource.data.uid == request.auth.uid;
      allow update, delete: if request.auth != null && resource.data.uid == request.auth.uid;
    }
    
    // Regole per la collezione dei preferiti
    match /favorites/{favoriteId} {
      allow read, write: if request.auth != null && resource.data.userId == request.auth.uid;
    }
    
    // Regole per la collezione notifiche
    match /notifications/{notificationId} {
      allow read, update, delete: if request.auth != null && resource.data.recipientId == request.auth.uid;
      allow create: if request.auth != null;
    }
  }
}
\end{lstlisting}
\end{scriptsize}

\clearpage

\newchapter{Progettazione e sviluppo applicazione in Kotlin}
\newsection{Architettura del sistema}

L'applicazione Android segue le linee guida di Google per l'architettura delle applicazioni moderne (Modern Android Architecture) basandosi sui principi della \textbf{Clean Architecture} e sul concetto fondamentale di \textbf{"Separation of Concerns"} (separazione delle responsabilità). Ciò consente di garantire un'elevata, manutenibilità e facilità di estensione delle componenti:

\begin{itemize}
    \item \textbf{Data Layer}: gestisce l'accesso, la persistenza e la sincronizzazione remota dei dati. Interagisce con i servizi cloud attraverso l'integrazione di \textbf{Firebase Cloud Firestore} utilizzato per la sincronizzazione asincrona di annunci, notifiche e preferiti, \textbf{Firebase Authentication} per la gestione sicura delle sessioni utente, e \textbf{Supabase Client} per il caricamento efficiente delle immagini. Include inoltre la libreria \textbf{Glide} per il caching e la resa prestazionale dei contenuti multimediali.
    
    \item \textbf{Domain Layer}: definisce le entità di business pure del sistema (le classi modello come \texttt{User}, \texttt{PuppyPost} e \texttt{Notification}) e racchiude la logica applicativa indipendente dalle librerie esterne. Regola le operazioni principali del dominio, quali il filtraggio degli annunci e la gestione dei preferiti.
    
    \item \textbf{Presentation Layer}: contiene l'interfaccia utente visiva organizzata in \textit{Activities} (\texttt{HomeActivity}, \texttt{SignInActivity}, \texttt{SignUpActivity}) e \textit{Fragments} (\texttt{HomeFragment}, \texttt{PuppyDetailsFragment}, \texttt{ProfileFragment}, \texttt{FavoritesFragment}, \texttt{AddPuppyFragment}). L'interfaccia si avvale di componenti grafici dichiarati in layout XML e di \textit{RecyclerView} affiancate da \textit{Custom Adapters} (\texttt{FeedAdapter}, \texttt{SearchAdapter}, \texttt{FavoritesAdapter}). Gestisce l'interattività dell'utente e la risposta reattiva alle gesture native, come lo \textit{Swipe-to-Favorite} implementato tramite \texttt{ItemTouchHelper}.
\end{itemize}

\begin{figure}[htbp]
    \centering
    \begin{tikzpicture}[
        node distance=1.4cm,
        box/.style={draw=black, thick, rectangle, fill=white, minimum width=10.5cm, minimum height=1.1cm, align=center, font=\bfseries},
        labeltext/.style={font=\small, text=black},
        arrow/.style={<->, >=stealth, thick, color=black}
    ]

        % Nodi (Livelli)
        \node (data) [box] {Data Layer (Firebase Auth, Firestore DB \& Supabase Storage)};
        \node (domain) [box, below=of data] {Domain Layer (Models \& Business Logic)};
        \node (presentation) [box, below=of domain] {Presentation Layer (Activities, Fragments \& Custom Adapters)};

        % Frecce e Testi tra Data e Domain
        \draw[arrow] (data.south) -- (domain.north);
        \node[labeltext, left=0.2cm of $(data.south)!0.5!(domain.north)$] {Restituisce dati / Stream};
        \node[labeltext, right=0.2cm of $(data.south)!0.5!(domain.north)$] {Richiede dati / Chiamate Cloud};

        % Frecce e Testi tra Domain e Presentation
        \draw[arrow] (domain.south) -- (presentation.north);
        \node[labeltext, left=0.2cm of $(domain.south)!0.5!(presentation.north)$] {Notifica cambiamenti / Aggiorna UI};
        \node[labeltext, right=0.2cm of $(domain.south)!0.5!(presentation.north)$] {Richiede azioni / Invia filtri};

    \end{tikzpicture}
    \caption{Architettura a tre livelli dell'applicazione Kotlin Paws}
    \label{fig:architettura_paws}
\end{figure}

\clearpage
\newsubsection{Mappa globale dell'architettura}

Il diagramma riportato in Figura~\ref{fig:mappa_globale_architettura} illustra la mappa globale dell'architettura dell'applicazione Android \textbf{Paws}, evidenziando l'organizzazione strutturale dei moduli e le relazioni di dipendenza tra le diverse componenti del sistema:

\begin{figure}[htbp]
    \centering
    \includegraphics[width=1\textwidth]{MappaGlobale.png}
    \caption{Diagramma delle componenti e mappa delle dipendenze dell'architettura Android Paws}
    \label{fig:mappa_globale_architettura}
\end{figure}


\vspace{0.5cm}
\newsubsection{Architettura d'avvio e navigazione}

L'avvio dell'applicazione è guidato dalla \texttt{SplashActivity}, la quale verifica lo stato dell'autenticazione dell'utente tramite il client di \textbf{Firebase Authentication}:

\begin{itemize}
    \item Se non viene rilevata alcuna sessione utente attiva (\texttt{currentUser == null}), il flusso viene reindirizzato alla schermata \texttt{SignInActivity}, dalla quale l'utente può eventualmente accedere alla schermata di registrazione \texttt{SignUpActivity}.
    
    \item Se la sessione è attiva (oppure a seguito del login/registrazione completati con successo), l'utente accede alla schermata principale \texttt{HomeActivity}.
    
    \item La \texttt{HomeActivity} ospita il contenitore principale dei contenuti ed abilita la navigazione tra le 4 schede della barra inferiore (\textbf{BottomNavigationView}):
    \begin{enumerate}
        \item \textbf{Home} (\texttt{HomeFragment}): mostra il feed globale dei cuccioli, la barra di ricerca utenti.
        \item \textbf{Aggiungi Cucciolo} (\texttt{AddPuppyFragment}): permette la selezione dell'immagine dalla galleria nativa e l'apertura della sezione di creazione post (\texttt{CreatePostFragment}).
        \item \textbf{Preferiti} (\texttt{FavoritesFragment}): raccoglie l'elenco dei cuccioli salvati dall'utente per la consultazione rapida.
        \item \textbf{Profilo} (\texttt{ProfileFragment}): mostra i dettagli del profilo personale e la griglia dei post pubblicati dall'utente corrente.
    \end{enumerate}
\end{itemize}

\begin{figure}[htbp]
    \centering
    \includegraphics[width=1\textwidth]{AvvioENavigazione.drawio.png}
    \caption{Mappa dell'architettura d'avvio e del flusso di navigazione dell'applicazione Paws}
    \label{fig:avvio_navigazione_paws}
\end{figure}
\clearpage

\vspace{0.5cm}
\newsubsection{Architettura gestione utente e autenticazione}

Le funzionalità relative all'identità utente, alla registrazione e alla personalizzazione del profilo sono strutturate secondo il seguente schema architetturale:

\begin{itemize}
    \item \textbf{View}: le schermate \texttt{SignInActivity}, \texttt{SignUpActivity}, \texttt{ProfileFragment} e \texttt{UserProfileFragment} gestiscono l'input utente e la resa dell'interfaccia visiva.
    \item \textbf{Handler}: \texttt{ProfileImageManager} gestisce la selezione dell'immagine dell'utente prima dell'upload.
    \item \textbf{Repository}: \texttt{AuthRepository} gestisce la creazione dell'account e l'autenticazione remota su Firebase Auth, mentre \texttt{SupabaseManager} si occupa del caricamento sicuro dell'immagine profilo nel bucket di storage remoto.
    \item \textbf{Model \& Database}: l'entità \texttt{User} modella i dati dell'utente, sincronizzati in tempo reale su \textbf{Cloud Firestore}.
\end{itemize}

\begin{figure}[htbp]
    \centering
    \includegraphics[width=1\textwidth]{GestioneUtente.drawio.png}
    \caption{Architettura e flusso dati per la Gestione Utente e Autenticazione dell'applicazione Paws}
    \label{fig:gestione_utente_paws}
\end{figure}
\clearpage

\vspace{0.5cm}
\newsubsection{Architettura gestione cuccioli e annunci}

La gestione degli annunci dei cuccioli rappresenta il cuore dell'applicazione e sfrutta i seguenti moduli dedicati:

\begin{itemize}
    \item \textbf{View}: \texttt{HomeFragment}, \texttt{AddPuppyFragment}, \texttt{CreatePostFragment} e \texttt{PuppyDetailsFragment} osservano i dati e reagiscono agli eventi dell'utente.
    \item \textbf{Adapters \& Handlers}: \texttt{FeedAdapter}, \texttt{SearchAdapter}, \texttt{PostAdapter} e \texttt{GridPostAdapter} popolano le \textit{RecyclerView}. L'\texttt{ItemTouchHandler} intercetta le gesture di \textit{swipe-to-favorite} per aggiungere istantaneamente un cucciolo ai preferiti dal feed.
    \item \textbf{Repository \& Services}: \texttt{PostRepository} esegue le query NoSQL su \textbf{Cloud Firestore} per il recupero e filtraggio degli annunci, mentre \texttt{SupabaseManager} memorizza le foto dei cuccioli.
    \item \textbf{Model}: \texttt{PuppyPost} incapsula le informazioni del cucciolo (nome, razza, età, sesso, posizione e riferimento al proprietario).
\end{itemize}

\begin{figure}[htbp]
    \centering
    \includegraphics[width=1.0\textwidth]{GestioneCucciolo.drawio.png}
    \caption{Architettura e flusso dati per la Gestione Cuccioli e Annunci dell'applicazione Paws}
    \label{fig:gestione_cucciolo_paws}
\end{figure}
\clearpage

\vspace{0.5cm}
\newsubsection{Architettura gestione notifiche e preferiti}

La gestione dei preferiti e delle notifiche garantisce un'esperienza d'uso reattiva e sincronizzata:

\begin{itemize}
    \item \textbf{View}: \texttt{FavoritesFragment} e \texttt{NotificationsFragment} mostrano rispettivamente la lista dei cuccioli salvati e le interazioni/notifiche ricevute.
    \item \textbf{Adapter}: \texttt{FavoritesAdapter} e \texttt{NotificationAdapter} formattano la resa grafica degli elementi. La libreria esterna \textbf{Glide} gestisce il caricamento fluido e asincrono delle immagini con caching in memoria.
    \item \textbf{Model}: \texttt{Notification} modella le notifiche di sistema salvate su Firestore.
\end{itemize}

\begin{figure}[htbp]
    \centering
    \includegraphics[width=1\textwidth]{GestioneNotifiche.drawio.png}
    \caption{Architettura e flusso dati per la Gestione Cuccioli e Annunci dell'applicazione Paws}
    \label{fig:gestione_notifiche_paws}
\end{figure}

\clearpage
\newsection{Mockup}
Il mockup dell’applicativo GESTRi è stato realizzato con Figma, uno strumento per la progettazione e la prototipazione di interfacce utente. Ogni mockup rappresenta la versione progettuale dell’interfaccia, definendo struttura, elementi grafici e funzionalità prima della fase di sviluppo.



\clearpage
\newsection{Verifica e Validazione del Software: Unit Tests}
I seguenti test rappresentano l’insieme delle attività di verifica implementate per garantire il corretto funzionamento del progetto. Sono stati adottati due approcci complementari: test automatici lato codice e test manuali delle API.

\newsubsection{Test Manuali} 
Per verificare le API è stata utilizzata l’applicazione Postman, uno strumento che consente di inviare richieste HTTP e analizzare le risposte. Questi test non rientrano nel diagramma degli Unit Test, ma rappresentano un'importante componente di validazione dell’integrazione e del comportamento delle API lato client-server.

\newsubsection{Unit Test} 
La sezione presenta i test unitari eseguiti sul modulo Attività di GESTRi, con l’obiettivo di verificarne la correttezza funzionale e assicurare l’affidabilità del relativo codice.
\begin{scriptsize}
\estiloPython
\begin{lstlisting}[caption={Attivita}, label=lst:javacode]
class AttivitaIntegrationTests(TestCase):
   def setUp(self):
       self.client = Client()
       # create users
       self.staff = Utente.objects.create(email='staff@example.com', ruolo=Ruolo.STAFF)
       self.staff.set_password('pass')
       self.staff.save()

\end{lstlisting}
\end{scriptsize}

\end{document}
